% A.64: root and PBC full proof/source/certificate audits PASS.
\subsection{A prime-to-closure enlargement criterion and all Heisenberg enlargements}

\begin{statement}{A.64}{Proved}{Torsion enlargement outside the geometric closure order}
Let \(Y\to U=E\setminus\{O\}\) be any geometrically
connected finite etale cover from Task A.5, with degree
\(d\), proper compactification \(C\), genus \(g\), and
extended map \(q:C\to E\). Let \(N\) be the order of
its geometric Galois closure. If an integer \(m\ge1\)
has a prime divisor \(r\nmid N\), then
\[
 Y_m=C\setminus q^{-1}(E[m])
                \xrightarrow{[m]q}U                   \tag{A.64.1}
\]
has full pro-\(\ell\) fixed group equal to the selected
cusp inertia, for every coefficient prime \(\ell\),
every cusp packet, and every permitted open derived
subgroup. It has degree \(m^2d\), proper genus \(g\),
and \(m^2d-2g+2\) cusps. The residue prime \(r=2\)
is allowed. If \(Y\) has pro-\(\ell\) cuspidal
normalizer recognition, then so does \(Y_m\).

As a consequence, for every odd prime \(p\) and every
integer \(m\ge1\), both assertions of Task A.5 hold
for the actual cover
\[
 C_p\setminus q_p^{-1}(E[pm])
                  \xrightarrow{[pm]q_p}U,             \tag{A.64.2}
\]
where \(C_p\) is the smooth proper normalization of
\(z^p=\psi_p(x)\) and \(q_p:C_p\to E\).
Its degree is \(p^3m^2\), genus
\(1+p^2(p-1)/2\), and cusp count
\(p^3m^2-p^3+p^2\).

The map \(Y_m\hookrightarrow Y\) is an open immersion.
These assertions do not give a finite refinement over the
original map \(Y\to U\), and they do not lift fixed
elements from a quotient fundamental group. The universal
Task A.5 remains open.
\end{statement}

\begin{proof}
\paragraph{An actual good tame model at every allowed residue prime.}
Choose a place above \(r\). The elliptic curve
\(E:y^2=x^3-x\) has potentially good reduction there
because \(j(E)=1728\) is integral; this includes
\(r=2\), by \cite[Proposition VII.5.5, p.197,
and Appendix A.1.4(b)]{fghSilvermanAEC}.
After finite extension, take its smooth proper group
scheme with the smooth origin section.

Since the actual geometric Galois closure has order
prime to \(r\), prime-to-\(r\) specialization for
the good pair \((E,O)\) extends that cover tamely
over this model after finite extension. The relevant
comparison, including inertia, is
\cite[Expos\'e XIII, Corollary 2.9, Section 2.10 and
the proof of Corollary 2.12]{fghSGA1}. This is the
same actual finite-cover comparison proved in A.56
and used in A.59: a geometric prime-to-\(r\)
cover of the generic fibre corresponds to a tame
cover of the special fibre, and finite descent
defines the required cover and its generic
identification over a finite extension. A possible
arithmetic twist is removed by that extension.

Normalize the good elliptic model in the actual
function field of \(C\). The resulting proper
source is smooth. Away from the origin it is
etale. At a point above the origin, after strict
henselian localization, the tame local map is
\[
                    R[[t]]\longrightarrow R[[u]],
                              t=u^e,\quad r\nmid e.
                                                        \tag{A.64.3}
\]
The indices divide the prime-to-\(r\) closure
order, so this argument applies to subgroup
quotients as well as to the Galois closure.
Thus the conclusion is potentially good
\emph{curve} reduction of \(C\), with the actual
map finite and etale outside the origin, rather
than good reduction of its Jacobian alone.
Nothing in this argument excludes residue two.

Enlarge the field further to split every point
in the finite set \(q^{-1}(E[m])\). The good
source and target and the tame local descriptions
persist under this base change. A nonconstant map
to an elliptic curve forces \(g\ge1\).

\paragraph{Every target torsion point has a nontrivial specialization cluster.}
Write \(m=r^a n\), where \(a\ge1\) and \(r\nmid n\).
Let \(f\in\{0,1\}\) be the \(r\)-rank of the good
elliptic special fibre. The geometric special
\(r^a\)-torsion has \(r^{af}\) points, while its
generic fibre has \(r^{2a}\) points. These facts
hold also in characteristic two; see
\cite[Theorem V.3.1, p.145]{fghSilvermanAEC}.

For precision, every special point lifts to a
generic torsion point after finite extension,
and every specialization fibre has the same size.
Indeed the finite flat group scheme \(E[r^a]\)
over a strict henselian discrete valuation ring
is the disjoint union of the finite flat pieces
whose special fibres are its individual geometric
points: lift the idempotents of its special
coordinate algebra. Each piece has a nonempty
generic fibre. Translation by a generic lift of
a special point permutes these pieces, so all
pieces have the same rank. Their total rank is
\(r^{2a}\). The generic fibre is etale, and hence
each piece contains exactly
\[
                             c=r^{a(2-f)}\ge2           \tag{A.64.4}
\]
generic geometric points. The \(n\)-torsion is
finite etale and specializes bijectively after
splitting. Decomposing \(E[m]\) into its
\(r\)-primary and prime-to-\(r\) parts therefore
proves that \emph{every} geometric specialization
cluster of \(E[m]\) has exactly \(c\) points.
This includes clusters away from the origin when
the reduction is ordinary; they are not omitted.

\paragraph{The exact source clusters.}
Let \(b\) be a special point of the target
underlying one of these torsion clusters.
If \(b\ne O\), the finite source map is etale
above its whole formal disc. Each special source
point above \(b\) receives precisely one lift
of each of the \(c\) generic target points,
so its source cluster has size \(c\ge2\).

If \(b=O\), let \(P\) be an original cusp
with index \(e_P\). The target cluster contains
the origin itself and \(c-1\) other points.
Equation (A.64.3) gives one point above the
origin and \(e_P\) distinct characteristic-zero
points above each other target point. All of
them specialize to \(P\). The source cluster
size is exactly
\[
                              1+e_P(c-1)\ge2.          \tag{A.64.5}
\]
The unit extraction used in (A.64.3) is legitimate
because \(e_P\) is invertible in the residue
field. These two cases account for every mark
of \(q^{-1}(E[m])\).

The proper source model is smooth. Distinct
marked sections specializing to one of its
points separate, after semistable marked
resolution, on a rational tree attached at
that point. Since every cluster has at least
two marks, every selected cusp has a distinct
companion whose entire joining path, including
the endpoints, is rational. This is the exact
pointed hypothesis of A.41--A.42. It yields the
full integral fixed-inertia conclusion at
every coefficient prime, including \(\ell=r\),
and for every packet. The finite global field
restriction can be chosen with its closed
derived subgroup inside the prescribed open
derived subgroup, exactly as in A.41.

The original indices satisfy
\(\sum_P e_P=d\) and
\(2g-2=\sum_P(e_P-1)=d-|q^{-1}(O)|\).
In characteristic zero, all other target
points are unramified. Hence
\[
 \begin{split}
 |q^{-1}(E[m])|
   &=|q^{-1}(O)|+d(m^2-1)\\
   &=m^2d-2g+2.
 \end{split}                                            \tag{A.64.6}
\]
The map (A.64.1) is finite etale of degree
\(m^2d\), and its source remains geometrically
connected. The normalizer assertion is the
actual-open-immersion argument of A.53 using
A.11 and unchanged proper homology; it does
not require a fixed-element lifting assertion.

\paragraph{Every odd Heisenberg prime and every enlargement.}
Let \(Y_p=C_p\setminus q_p^{-1}(E[p])\).
Its degree-\(p^3\) map to \(U\) is
\(f_p=[p]q_p\), whose geometric Galois closure
has the actual order \(p^3\), by A.61--A.62.
Moreover
\[
                   f_p^{-1}(E[m])=q_p^{-1}(E[pm]).       \tag{A.64.7}
\]
If \(m\) has a prime factor different from
\(p\), apply the general conclusion above to
\(f_p\). This treats every such \(m\), with
no congruence restriction on that other prime.
If \(p\equiv1\pmod4\), A.63 treats all
remaining \(m\) as well.

It remains to take \(p\equiv3\pmod4\) and
\(m=p^a\), \(a\ge0\). Use the actual
compact-type model of A.61 at residue \(p\).
The target elliptic curve is supersingular,
so all points of \(E[p^{a+1}]\) specialize to
\(O\). Its proper source has one elliptic
main component mapping radicially to the
good target, and its other positive-genus
components map constantly to the finite
zeros of the function \(R\) in A.61.
None of these centres is \(O\), where
\(R\) has its pole.

Consequently no enlarged mark can specialize
on any of those positive-genus tails. On the
main component there is exactly one point
mapping to \(O\), by radiciality. Every
rational component carrying an enlarged
mark lies in the fibre over \(O\): its map
to a good elliptic curve is constant. The
path from such a component to the main
cannot pass through a positive-genus tail,
because the map at each connecting node
would then have both values \(O\) and a
different critical centre. All rational
parts of this fibre attach at the unique
main point over \(O\). Nodalness allows
only one branch at that point on the main,
so these parts form one connected rational
subtree, if they are present.

If there is no such rational part in the
unmarked model, all marks specialize to
that single smooth main point and marked
separation creates a rational tree there.
If there is already a rational part, the
main attachment is a node, so the smooth
marked sections lie on that rational part
or on rational components inserted above
its nodes. Thus every joining path between
the enlarged marks is rational. The set
contains the original \(p^2\ge9\) cusps,
so every mark has a companion. A.41--A.42
again prove all-prime fixedness.

A.61--A.62 give all-prime normalizer
recognition on \(Y_p\); A.53's actual open
immersion argument gives it on every
enlargement (A.64.2). This completes the
asserted consequence while retaining the
universal open frontier stated above.

The executed certificate
\path{code/a_torsion_enlargement_closure.py}
checks both specialization-cluster counts
and the Euler identity for actual
Heisenberg degrees at several primes and
enlargements, including residue two.
\end{proof}
